\section{Conclusion}
\label{sec:conclusion}

We investigated whether conditional diffusion augmentation can generate physically
plausible extreme coastal water-level residuals that improve downstream forecasting
on independent real storm events. The answer is \emph{not yet}: our pilot DDPM
produced samples whose standard deviation is approximately 23 times the training
target SD, with magnitudes exceeding $\pm 16$ metres. This failure is reproduced
on both CUDA and CPU, and traced to terminal-step amplification in the cosine
noise schedule combined with imperfect epsilon prediction at a limited training
budget.

More importantly, this work establishes \emph{what is needed} before diffusion
augmentation can be meaningfully evaluated for coastal water-level forecasting:
\begin{enumerate}
    \item A frozen data protocol with signed-off licenses, datums, and canonical
    storm event definitions.
    \item Strong neural forecasting baselines (GRU/TCN) on the formal 96h$\rightarrow$72h
    release contract.
    \item A controlled retraining of the diffusion model with investigation of
    schedule, objective, and training budget.
    \item Downstream forecasting evaluation on independent real storm events with
    event-level paired bootstrap confidence intervals.
\end{enumerate}

All of this infrastructure is now in place. The \texttt{trae/research-recovery-20261030}
branch contains the merged PR \#19 (data audit, 55 tests) and PR \#21 (diffusion
implementation, 19 tests) code, with 69 tests passing, zero Ruff warnings, and
a repaired embargo filter. The frozen data protocol, embargo-safe splitting, and
engineering validation guards prevent reintroduction of known leakage and
generation-quality failures in future work.

The claim in \texttt{paper/CLAIMS\_LEDGER.md} remains exactly as it was at the
start: H1 and H2 are \textbf{untested hypotheses}. The contribution of this paper
is to record honestly what we have found, what we have ruled out, and what
remains to be done.
